% ---------------------------------------------------------------------------
%  V. Recovered Command Surface and Edge Platform
%  Floats:  none.
%  Rule:    concepts, not identifiers. No source filenames, no method names,
%           no vendor package names. The one deliberate exception is the
%           broadcast listing below, which IS the contribution of C2.
%           This is a standing reviewer request; do not reintroduce them.
% ---------------------------------------------------------------------------

\section{Recovered Command Surface and Edge Platform}
\label{sec:platform}

\subsection{Command and telemetry}

Motion is commanded by advertising a multiplexer input and then publishing twist messages on it
at roughly \SI{10}{\hertz}. The vendor application quietly sets a manual-control mode when a
session opens. Our early builds did not, and every command we sent was dropped without comment.
That precondition was visible in the decompiled application and in no traffic capture we took,
and it was our first concrete encounter with H4. Autonomous motion comes in two forms: a
coordinate goal, and a named-annotation service backed by the vendor's marker manager. One
status field reports the state machine gating both, with values for unlocalized, ready,
navigating, arrived and error. We throttle subscriptions per topic to keep the bridge usable
over the robot's Wi-Fi.

One detail cost us time and is worth passing on. A single vendor message type, carrying one
integer field, is reused across several unrelated services, each giving that integer its own
private meaning. Call one of them with empty or guessed arguments and it returns success, then
either does nothing or does the opposite of what was intended. We recovered the correct
arguments through the same introspection route that supports H1. Interfaces reconstructed from
the outside are prone to this, and an acknowledgment offers no protection against it.

\subsection{Edge architecture}

The reconstructed platform is edge-first by necessity, since the robot's network offers no
guaranteed Internet path. A domain layer hides the robot behind one interface with two
implementations, simulated and real. This let us exercise the logic under 169 unit tests without
the physical machine, and it kept ordinary software defects separate from genuine protocol
uncertainty, which during field sessions was the difference between a productive afternoon and a
wasted one. Above it, an application layer exposes REST and WebSocket endpoints and runs either
on a development machine or on the Android head, so a visitor session needs no external
computer. A presentation layer serves the visitor interface of Fig.~\ref{fig:robot} and the
operator console.

\subsection{Speech through inter-process communication}

Speech was the hardest capability to recover, and at the protocol level the search simply
failed. No topic accepted text, no such service existed, the head's HTTP ports were closed, and
the broker carried telemetry only. The reason turned out to be architectural rather than
accidental. The vendor implements speech inside the Android application layer, over opaque IPC
to a bundled engine, so it was never a robot-level capability at all and no amount of protocol
search would have found it. The answer was to stop looking sideways and move up a layer. A
\SI{50}{\kilo\byte} companion application registers a broadcast
receiver~\cite{chin2011android} and hands incoming text to the platform's native speech engine:

\begin{footnotesize}
\begin{verbatim}
am broadcast -n com.cybel.ttsbridge/.SpeakReceiver
  -a com.cybel.ttsbridge.SPEAK --es text "Welcome."
\end{verbatim}
\end{footnotesize}

\noindent The move itself generalises. Whenever a vendor implements a function in the
application layer, no protocol-level search will reach it, and the host operating system is
where the capability has to be picked up instead.

\subsection{Interaction layer}

Visitor interaction runs entirely on the head. We abandoned free dictation early, because a
compact offline model has no pronunciation for the proper nouns that matter on a campus;
recognition now works from a closed vocabulary rebuilt at request time from the deployed
annotations and the question set. The same limitation caught the wake phrase, built from the
platform's invented name and consistently decoded as an unrelated word the model did know, until
we replaced it with a phonetically close phrase of real words. Face re-identification runs at
roughly \SI{1}{\hertz} and sends embedding vectors only, never images.
